\documentclass[12pt, a4paper]{article}
\usepackage[T2A]{fontenc}
\usepackage[utf8]{inputenc}
\usepackage[russian]{babel}
\usepackage{amsmath, amssymb}
\usepackage{enumitem}
\usepackage{geometry}
\geometry{top=2cm, bottom=2cm, left=2.5cm, right=2cm}

\begin{document}

\section*{Георгий Игоревич Вольфсон <<Делимость с человеческим лицом>>}

\subsection*{Беседа седьмая. Остатки.}

\subsubsection*{Упражнения}

\begin{enumerate}[label=\arabic*., wide=0pt, leftmargin=*]

\item \textbf{Найдите остаток от деления \(2019 - 2020 \cdot 2021\) на 7.}

\textbf{Решение:}
Найдём остатки чисел при делении на 7:
\[
2019 \equiv 3 \pmod{7},\quad 2020 \equiv 4 \pmod{7},\quad 2021 \equiv 5 \pmod{7}.
\]
Тогда:
\[
2019 - 2020 \cdot 2021 \equiv 3 - 4 \cdot 5 = 3 - 20 = -17 \pmod{7}.
\]
Так как \(-17 = -21 + 4 = 7 \cdot (-3) + 4\), то остаток равен 4.
Ответ: 4.

\item \textbf{Найдите все натуральные числа, у которых остаток от деления на 11 вдвое больше неполного частного.}

\textbf{Решение:}
Пусть \(a\) — искомое число, \(q\) — неполное частное. Тогда:
\[
a = 11q + 2q = 13q,
\]
причём остаток \(2q\) должен быть меньше делителя 11:
\[
2q < 11 \Rightarrow q = 1,2,3,4,5.
\]
Получаем числа:
\[
q=1 \Rightarrow 13,\quad q=2 \Rightarrow 26,\quad q=3 \Rightarrow 39,\quad q=4 \Rightarrow 52,\quad q=5 \Rightarrow 65.
\]
Ответ: 13, 26, 39, 52, 65.

\item \textbf{Найдите остаток от деления \(12^{100}\) на 11.}

\textbf{Решение:}
Так как \(12 \equiv 1 \pmod{11}\), то:
\[
12^{100} \equiv 1^{100} = 1 \pmod{11}.
\]
Ответ: 1.

\item \textbf{Найдите остаток от деления \(13^{2019}\) на 7.}

\textbf{Решение:}
Так как \(13 \equiv 6 \pmod{7}\), то:
\[
13^{2019} \equiv 6^{2019} \pmod{7}.
\]
Заметим, что \(6 \equiv -1 \pmod{7}\), значит:
\[
6^{2019} \equiv (-1)^{2019} = -1 \equiv 6 \pmod{7}.
\]
Ответ: 6.

\item \textbf{Найдите последнюю цифру числа \(17^{15} + 15^{17}\).}

\textbf{Решение:}
Последняя цифра \(17^{15}\) определяется последней цифрой основания 7. Цикл последних цифр степеней 7: \(7, 9, 3, 1\) (период 4). Так как \(15 \equiv 3 \pmod{4}\), то последняя цифра \(17^{15}\) равна 3.
Число \(15^{17}\) оканчивается на 5 (так как любая степень числа, оканчивающегося на 5, оканчивается на 5).
Сумма: \(3 + 5 = 8\).
Ответ: 8.

\item \textbf{Найдите последнюю цифру числа \(18^{18^{18}}\).}

\textbf{Решение:}
Последняя цифра \(18^k\) определяется последней цифрой 8. Цикл последних цифр степеней 8: \(8, 4, 2, 6\) (период 4). Нужно найти \(18^{18} \bmod 4\).
Так как \(18 \equiv 2 \pmod{4}\), то:
\[
18^{18} \equiv 2^{18} \pmod{4}.
\]
При \(k \ge 2\) число \(2^k\) делится на 4, поэтому \(2^{18} \equiv 0 \pmod{4}\).
Значит, \(18^{18}\) при делении на 4 даёт остаток 0, то есть показатель степени кратен 4. Тогда последняя цифра \(18^{18^{18}}\) равна последней цифре \(8^4\), то есть 6.
Ответ: 6.

\item \textbf{Найдите последнюю цифру числа \(1^5 + 2^5 + \dots + 100^5\).}

\textbf{Решение:}
Легко проверить, что для любого \(n\) последняя цифра \(n^5\) совпадает с последней цифрой \(n\) (это следует из малой теоремы Ферма или прямой проверки для цифр 0–9). Тогда:
\[
1^5 + 2^5 + \dots + 100^5 \equiv 1 + 2 + \dots + 100 \pmod{10}.
\]
Сумма \(1 + 2 + \dots + 100 = \frac{100 \cdot 101}{2} = 5050\), последняя цифра 0.
Ответ: 0.

\item \textbf{Решите уравнение в целых числах: \(x^2 + 24xy = 2021\).}

\textbf{Решение:}
Рассмотрим уравнение по модулю 3:
\[
24xy \equiv 0 \pmod{3},\quad 2021 \equiv 2 \pmod{3}.
\]
Тогда \(x^2 \equiv 2 \pmod{3}\). Но квадрат целого числа по модулю 3 может давать только остатки 0 или 1 (так как \(0^2=0\), \(1^2=1\), \(2^2=4 \equiv 1\)). Противоречие.
Ответ: решений нет.

\item \textbf{Решите уравнение в целых числах: \(x^2 - 16xy + y^2 = 2019\).}

\textbf{Решение:}
Рассмотрим уравнение по модулю 4:
\[
16xy \equiv 0 \pmod{4},\quad 2019 \equiv 3 \pmod{4}.
\]
Тогда \(x^2 + y^2 \equiv 3 \pmod{4}\). Но квадрат целого числа по модулю 4 может давать только остатки 0 или 1. Следовательно, \(x^2 + y^2\) может давать остатки \(0, 1\) или \(2\), но не 3. Противоречие.
Ответ: решений нет.

\item \textbf{Решите уравнение в целых числах: \(x^3 + 9x^2y = 2020\).}

\textbf{Решение:}
Рассмотрим уравнение по модулю 9:
\[
9x^2y \equiv 0 \pmod{9},\quad 2020 \equiv 4 \pmod{9}.
\]
Тогда \(x^3 \equiv 4 \pmod{9}\). Но куб целого числа по модулю 9 может давать только остатки 0, 1 или 8 (проверка: \(0^3=0\), \(1^3=1\), \(2^3=8\), \(3^3=27\equiv0\), \(4^3=64\equiv1\), \(5^3=125\equiv8\), \(6^3=216\equiv0\), \(7^3=343\equiv1\), \(8^3=512\equiv8\)). Остаток 4 невозможен.
Ответ: решений нет.

\item \textbf{Найдите все простые числа \(p\), для которых \(p^2 + 20\) также простое число.}

\textbf{Решение:}
Рассмотрим \(p\) по модулю 3. Если \(p \neq 3\), то \(p \equiv \pm 1 \pmod{3}\), значит \(p^2 \equiv 1 \pmod{3}\). Тогда:
\[
p^2 + 20 \equiv 1 + 2 \equiv 0 \pmod{3}.
\]
То есть \(p^2 + 20\) делится на 3. Так как \(p^2 + 20 > 3\), то оно составное.
Остаётся проверить \(p = 3\): \(3^2 + 20 = 9 + 20 = 29\) — простое число.
Ответ: \(p = 3\).

\item \textbf{Докажите, что \(n^5 + 4n\) делится на 5 при любом натуральном \(n\).}

\textbf{Решение:}
Составим таблицу остатков при делении \(n\) на 5:
\[
\begin{array}{c|ccccc}
n \bmod 5 & 0 & 1 & 2 & 3 & 4 \\
\hline
n^5 \bmod 5 & 0 & 1 & 2^5=32\equiv2 & 3^5=243\equiv3 & 4^5=1024\equiv4 \\
4n \bmod 5 & 0 & 4 & 8\equiv3 & 12\equiv2 & 16\equiv1 \\
n^5 + 4n \bmod 5 & 0 & 5\equiv0 & 5\equiv0 & 5\equiv0 & 5\equiv0
\end{array}
\]
Во всех случаях остаток равен 0, значит \(n^5 + 4n\) делится на 5.

\item \textbf{Докажите, что \(n^3 - n\) делится на 6 при любом натуральном \(n\).}

\textbf{Решение:}
Составим таблицу остатков при делении \(n\) на 6:
\[
\begin{array}{c|cccccc}
n \bmod 6 & 0 & 1 & 2 & 3 & 4 & 5 \\
\hline
n^3 \bmod 6 & 0 & 1 & 8\equiv2 & 27\equiv3 & 64\equiv4 & 125\equiv5 \\
n \bmod 6 & 0 & 1 & 2 & 3 & 4 & 5 \\
n^3 - n \bmod 6 & 0 & 0 & 0 & 0 & 0 & 0
\end{array}
\]
Во всех случаях остаток равен 0, значит \(n^3 - n\) делится на 6.

\item \textbf{Найдите все натуральные \(n\) от 50 до 70, для которых \(n^3 + 2n^2 + 7\) кратно 5.}

\textbf{Решение:}
Составим таблицу остатков при делении \(n\) на 5:
\[
\begin{array}{c|ccccc}
n \bmod 5 & 0 & 1 & 2 & 3 & 4 \\
\hline
n^3 \bmod 5 & 0 & 1 & 8\equiv3 & 27\equiv2 & 64\equiv4 \\
2n^2 \bmod 5 & 0 & 2 & 8\equiv3 & 18\equiv3 & 32\equiv2 \\
7 \bmod 5 & 2 & 2 & 2 & 2 & 2 \\
\hline
\text{Сумма} \bmod 5 & 2 & 5\equiv0 & 8\equiv3 & 7\equiv2 & 8\equiv3
\end{array}
\]
Выражение кратно 5 только при \(n \equiv 1 \pmod{5}\). В диапазоне от 50 до 70:
\[
51,\;56,\;61,\;66.
\]
Ответ: 51, 56, 61, 66.

\item \textbf{Натуральные числа \(x, y, z\) таковы, что \(x^2 + y^2 = z^2\). Докажите, что \(xyz\) делится на 60.}

\textbf{Решение:}
Нужно доказать, что \(xyz\) делится на 3, 4 и 5.

\textbf{Делимость на 3:} 
Квадрат числа по модулю 3 даёт остаток 0 или 1. Если ни одно из чисел \(x, y, z\) не делится на 3, то \(x^2 \equiv y^2 \equiv z^2 \equiv 1 \pmod{3}\), и равенство \(1 + 1 = 1\) невозможно. Значит, хотя бы одно из \(x, y, z\) делится на 3.

\textbf{Делимость на 4:}
Квадрат числа по модулю 4 даёт остаток 0 или 1. Если \(x\) и \(y\) нечётны, то \(x^2 \equiv y^2 \equiv 1 \pmod{4}\), значит \(z^2 \equiv 2 \pmod{4}\), что невозможно. Следовательно, хотя бы одно из чисел \(x, y\) чётно, значит \(xyz\) делится на 2. Для делимости на 4 нужно, чтобы одно из чисел делилось на 4 или оба были чётными. Если оба \(x, y\) чётны, то \(xy\) делится на 4. Если одно чётно, а другое нечётно, то \(z\) нечётно, а чётное число должно делиться на 4 (иначе квадрат чётного числа, не делящегося на 4, даёт остаток 4 при делении на 8, но нам достаточно, что если \(x\) или \(y\) не делится на 4, то другое должно делиться на 4 из соображений сравнения по модулю 4). В любом случае \(xyz\) делится на 4.

\textbf{Делимость на 5:}
Квадрат числа по модулю 5 даёт остатки 0, 1 или 4. Если ни одно из чисел \(x, y, z\) не делится на 5, то их квадраты могут быть только 1 или 4. Возможные суммы: \(1+1=2\), \(1+4=5\equiv0\), \(4+4=8\equiv3\). Ни одна из этих сумм не равна 1 или 4 (кроме случая \(1+4=0\), но тогда \(z^2 \equiv 0\), значит \(z\) делится на 5, противоречие). Значит, хотя бы одно из \(x, y, z\) делится на 5.

Таким образом, \(xyz\) делится на \(3 \cdot 4 \cdot 5 = 60\).

\end{enumerate}

\end{document}